\documentclass[10pt,letterpaper]{article}
\usepackage[T1]{fontenc}
\usepackage{mathptmx}
\usepackage[margin=0.6in,top=0.5in,bottom=0.6in]{geometry}
\usepackage{xcolor}
\usepackage{enumitem}
\usepackage{titlesec}
\usepackage{tabularx}
\usepackage{fancyhdr}
\usepackage{lastpage}
\usepackage[hidelinks,pdftitle={Wanli Qian -- Summer 2027 Research Internship},pdfauthor={Wanli Qian}]{hyperref}

\definecolor{navy}{HTML}{1F3A60}
\definecolor{todo}{HTML}{C0392B}
\hypersetup{colorlinks=true,urlcolor=navy,linkcolor=navy}

\pagestyle{fancy}
\fancyhf{}
\renewcommand{\headrulewidth}{0pt}
\fancyfoot[L]{\small\color{navy}Wanli Qian}
\fancyfoot[R]{\small\thepage\ / \pageref{LastPage}}

\setlength{\parindent}{0pt}
\titleformat{\section}{\color{navy}\bfseries}{}{0pt}{}[\vspace{-6pt}\color{black}\rule{\linewidth}{0.5pt}]
\titlespacing*{\section}{0pt}{9pt}{5pt}
\setlist[itemize]{leftmargin=12pt,itemsep=1pt,topsep=1pt,parsep=0pt,label=\textbullet,after=\vspace{3pt}}

\newcommand{\entry}[2]{\textbf{#1}\hfill\textit{#2}\\}
\newcommand{\sub}[2]{#1\hfill\textbf{#2}\\}
\newcommand{\pub}[3]{\begin{minipage}[t]{\linewidth}#1\\[-1pt]#2\\[-1pt]#3\end{minipage}\par\vspace{5pt}}
\newcommand{\me}{\textbf{Wanli Qian}}
\newcommand{\todo}[1]{\textcolor{todo}{[#1]}}

\begin{document}

{\fontsize{24}{26}\selectfont\bfseries Wanli Qian}\\[3pt]
\href{mailto:wanliqia@usc.edu}{wanliqia@usc.edu} \enspace|\enspace (678) 956-4378 \enspace|\enspace Los Angeles, CA\\
\href{https://silverwings-zero.github.io/}{Website} \enspace|\enspace
\href{https://www.linkedin.com/in/wanli-qian-7bb399192/}{LinkedIn} \enspace|\enspace
\href{https://github.com/silverwings-zero}{GitHub} \enspace|\enspace
\href{https://scholar.google.com/citations?user=azFoAIkAAAAJ&hl=en}{Google Scholar}\\[4pt]
\textbf{Seeking Summer 2027 Research Scientist Internship}\\
{\color{navy}Tactile State Estimation \,|\, Dexterous Manipulation \,|\, Teleoperation}

\section{Research Profile}
CS Ph.D. researcher studying how robots estimate the physical state and properties of objects through touch, multimodal
sensing, and action. Builds tactile-feedback teleoperation systems for contact-rich manipulation and data collection, and
develops learned models of force, vibration, and vision-based tactile signals. Research spans active perception under
uncertainty, manipulation, and trajectory optimization; recipient of a 2026 IEEE Haptics Symposium paper award.

\section{Education}
\entry{University of Southern California}{2024--Present}
\sub{Ph.D. in Computer Science}{Expected graduation: Fall 2028}
\entry{University of Chicago}{2022--2024}
Predoctoral M.S. in Computer Science\\
\entry{Georgia Institute of Technology}{2018--2022}
B.S. in Computer Science

\section{Technical Skills}
\textbf{Methods:} tactile perception (vision-based tactile, force, accelerometer), active perception, bilateral teleoperation,
haptic rendering, generative modeling, trajectory optimization, sampling- and graph-based motion planning, motion capture\\
\textbf{Tools:} \todo{List languages \& frameworks, e.g.\ Python, PyTorch, C++, ROS/ROS 2, GTSAM, MuJoCo/Isaac, GelSight/DIGIT --- keep only what you use}

\section{Research Experience}
\entry{USC HaRVI Lab \,|\, Ph.D. Researcher}{2024--Present}
\vspace{-11pt}
\begin{itemize}
  \item Convert vision-based tactile video of a grasped object into haptic feedback for the operator, conveying the contact
        interactions the gripper experiences during teleoperation (first-author RA-L submission, under review).
  \item Build bilateral teleoperation systems with tactile feedback for contact-rich manipulation, remote surface perception,
        and demonstration data collection.
  \item Develop active haptic perception policies that select exploratory robot actions to reduce uncertainty in estimates of
        object and surface physical properties.
  \item Curate multimodal contact datasets spanning 200+ materials, combining force, acceleration, audio, and vision-based
        tactile signals from tapping and sliding interactions.
  \item Develop language-conditioned generative models of tactile and vibration signals; integrate them into VR/AR
        haptic authoring interfaces (IEEE Haptics 2026, Distinguished Technical Long Paper Award).
\end{itemize}

\entry{University of Chicago AxLab \,|\, Predoctoral Researcher}{2022--2024}
\vspace{-11pt}
\begin{itemize}
  \item Led SHAPE-IT (UIST 2024, co-first author), an interactive system for conversational 3D shape generation and
        real-time reconfiguration of actuated physical shape displays.
  \item Developed adaptive LLM-based agents that modify shape-display control scripts in response to ongoing interaction.
\end{itemize}

\entry{Georgia Tech Borg Lab \,|\, Lab Assistant}{2021--2022}
\vspace{-11pt}
\begin{itemize}
  \item Developed motion planning and trajectory optimization methods for cable-driven parallel robots and articulated
        manipulators performing drawing and writing tasks (ICRA 2022).
  \item Integrated motion capture with robot control to track and reproduce human artistic trajectories; modeled strokes
        with polynomial fits and deformable virtual brush dynamics.
\end{itemize}

\section{Industry Experience}
\entry{Kolmostar \,|\, Software Engineer Intern}{2022}
\vspace{-11pt}
\begin{itemize}
  \item Built urban-geometry simulation pipelines to quantify wireless-signal shadowing, with visualization tools supporting
        sensing-system deployment.
\end{itemize}
\entry{Megvii \,|\, Robotics Software Solution Intern}{2019}
\vspace{-11pt}
\begin{itemize}
  \item Implemented and benchmarked probabilistic and graph-based motion planners for high-DOF industrial manipulators;
        analyzed efficiency, robustness, and convergence across task geometries.
\end{itemize}

\section{Publications}
{\color{navy}\textbf{Under Review}}\par\vspace{3pt}
\pub{\textbf{Feeling What the Gripper Feels: Haptic Rendering of a Grasped Object from Tactile Video in Teleoperation.}}
    {\me\ et al.}
    {Submitted to \textit{IEEE Robotics and Automation Letters (RA-L)}; under review.}
\pub{\textbf{TouchTwin: Human-in-the-Loop Haptic Texturing of 3D-Scanned Tabletop Scenes with Vision--Language Recognition and Language-Guided Generation.}}
    {\me\ et al.}
    {Submitted to \textit{ACM CHI 2027}; under review.}

{\color{navy}\textbf{Published}}\par\vspace{3pt}
\pub{\textbf{Language-Guided Multimodal Texture Authoring via Generative Models.}}
    {\me, Michael Gu, Aiden Chang, Shihan Lu, Heather Culbertson.}
    {\textit{IEEE Haptics Symposium}, 2026. \href{https://arxiv.org/abs/2604.06489}{Paper} | \href{https://youtu.be/f5vz52uPw-4}{Video} \enspace\textbf{Distinguished Technical Long Paper Award}}
\pub{\textbf{GTGraffiti: Spray Painting Graffiti Art from Human Painting Motions with a Cable Driven Parallel Robot.}}
    {Gerry Chen, Sereym Baek, JD Florez, \me, Sang-won Leigh, Seth Hutchinson, Frank Dellaert.}
    {\textit{IEEE ICRA}, 2022. \href{https://ieeexplore.ieee.org/document/9812008}{Paper} | \href{https://www.youtube.com/watch?v=R4ySYTNGv6s}{Video}}
\pub{\textbf{SHAPE-IT: Exploring Text-to-Shape-Display for Generative Shape-Changing Behaviors with LLMs.}}
    {\me*, Chenfeng Gao*, Anup Sathya, Ryo Suzuki, Ken Nakagaki.}
    {\textit{ACM UIST}, 2024. *Equal contribution. \href{https://dl.acm.org/doi/10.1145/3654777.3676348}{Paper} | \href{https://www.youtube.com/watch?v=2NxzOBc7AdQ}{Video}}
\pub{\textbf{Shape n'Swarm: Hands-on, Shape-aware Generative Authoring with Swarm UI and LLMs.}}
    {Matthew Jeung, Anup Sathya, \me, Steven Arellano, Luke Jimenez, Ken Nakagaki.}
    {\textit{ACM UIST}, 2025. \href{https://dl.acm.org/doi/10.1145/3746059.3747781}{Paper} | \href{https://www.youtube.com/watch?v=5u0M9yL7tyY}{Video}}
\pub{\textbf{Weakly Supervised Part-Based Method for Combined Object Detection in Remote Sensing Imagery.}}
    {\me, Zhe Yan, Zhe Zhu, Wei Yin.}
    {\textit{IEEE JSTARS}, vol.\ 15, pp.\ 5024--5036, 2022. \href{https://ieeexplore.ieee.org/document/9789416}{Paper}}

\section{Teaching}
\entry{CSCI 445: Introduction to Robotics \,|\, Teaching Assistant, USC}{2025}
\entry{CMSC 25040: Introduction to Computer Vision \,|\, Teaching Assistant, UChicago}{2023}

\section{Service \& Outreach}
\textbf{Reviewing:} CHI, UIST, IEEE Haptics Symposium, IEEE Transactions on Haptics.\\
\textbf{Mentorship:} USC CURVE (Spring/Fall 2025, Spring/Fall 2026); USC Viterbi SURE (Summer 2025); USC REU (Summer 2026).\\
\textbf{Outreach:} Peace4Kids robotics workshop organizer/mentor; RSS conference volunteer.

\end{document}
